%%
%% ACM SIGCONF Conference Paper Template for SciVLABench
%%

\documentclass[sigconf]{acmart}

%% BibTeX command to typeset BibTeX logo in the docs
\AtBeginDocument{%
  \providecommand\BibTeX{{%
    Bib\TeX}}}

%% Rights management information
\setcopyright{acmlicensed}
\copyrightyear{2025}
\acmYear{2025}
\acmDOI{XXXXXXX.XXXXXXX}

%% Conference information (TODO: Update with actual conference details)
\acmConference[Conference acronym 'XX]{Conference Title}{Month 00--00, 2025}{City, Country}

%%
%% The start of the document source.
\begin{document}

%%
%% Title
\title[SciVLABench]{SciVLABench: A Large-Scale Scientific Vision-Language-Action Benchmark}

%%
%% Authors (TODO: Update with real authors and affiliations)
\author{First Author}
\affiliation{
  \institution{Institution Name}
  \city{City}
  \state{State/Province}
  \country{Country}
}
\email{email@example.com}

\author{Second Author}
\affiliation{
  \institution{Institution Name}
  \city{City}
  \state{State/Province}
  \country{Country}
}
\email{email@example.com}

%% Short authors for page header
\renewcommand{\shortauthors}{Author et al.}

%%
%% Abstract
\begin{abstract}
Embodied agents are increasingly deployed in scientific laboratories, yet systematically benchmarking their capabilities remains an open challenge. While simulation-based evaluation has become the standard paradigm, existing benchmarks suffer from three key limitations. First, they lack a structured task framework, reducing evaluation to isolated operations rather than capturing higher-level competency dimensions. Second, manual task design fundamentally limits scale, failing to cover diverse instructions and complex, operation sequences. Third, the absence of external disturbances---such as human interventions---precludes the assessment of human--robot collaboration.

To fill this gap, we introduce SciVLABench, a large-scale, comprehensive benchmark for evaluating embodied scientific agents in realistic laboratory settings. In contrast to prior manual efforts, SciVLABench employs an automated task generation pipeline grounded in over 20 curated real-world experimental scenarios. It systematically evaluates eight core capability dimensions---including liquid transfer, physical mixing, thermodynamic control, equipment assembly, and experimental safety---through more than 1,000 complexity-stratified task templates, yielding a task pool 20 times larger than that of the previous largest benchmark. The tasks span basic actions, long-term autonomous experiments, and human--robot collaborative scenarios, providing a fine-grained testbed for model evaluation. SciVLABench is available at \url{https://github.com/kaikai-2022/SciVLABench}.
\end{abstract}

%%
%% CCS Concepts and Keywords (TODO: Generate proper CCS concepts from https://dl.acm.org/ccs/ccs.cfm)
\begin{CCSXML}
<ccs2012>
 <concept>
  <concept_id>00000000.0000000.0000000</concept_id>
  <concept_desc>Computing methodologies</concept_desc>
  <concept_significance>500</concept_significance>
 </concept>
 <concept>
  <concept_id>00000000.00000000.00000000</concept_id>
  <concept_desc>Artificial intelligence</concept_desc>
  <concept_significance>300</concept_significance>
 </concept>
 <concept>
  <concept_id>00000000.00000000.00000000</concept_id>
  <concept_desc>Robotics</concept_desc>
  <concept_significance>300</concept_significance>
 </concept>
</ccs2012>
\end{CCSXML}

\ccsdesc[500]{Computing methodologies}
\ccsdesc[300]{Artificial intelligence}
\ccsdesc[300]{Robotics}

%%
%% Keywords
\keywords{Scientific embodied agents, robotics benchmark, neuro-symbolic AI, vision-language-action models}

%%
%% Generate title page
\maketitle

%%
%% ==================== MAIN CONTENT ====================
\section{Introduction}

Driven by recent advances in multimodal perception, reasoning, and decision-making, artificial intelligence (AI) is rapidly extending from digital environments into the physical world through embodied agents. While remarkable progress has been achieved in domains such as household assistance and industrial automation, scientific laboratories remain one of the most challenging frontiers for embodied intelligence. Unlike general environments, laboratory tasks require agents to integrate dexterous manipulation, procedural understanding, scientific knowledge, and safety awareness throughout experimental execution. As embodied agents are increasingly envisioned as autonomous scientific assistants, systematically evaluating their readiness for laboratory environments has become an important and timely research problem.

Developing and evaluating embodied agents in real laboratories is prohibitively expensive, difficult to scale, and inherently risky. Consequently, high-fidelity simulation has emerged as an indispensable paradigm for developing and evaluating scientific embodied intelligence in a safe, scalable, and reproducible manner, leading to the emergence of simulation-based laboratory benchmarks. Nevertheless, a fundamental question remains largely unexplored: \textbf{how should scientific laboratory intelligence be systematically evaluated?}

Existing laboratory benchmarks primarily organize evaluation around execution characteristics, such as manipulation operations and task complexity (e.g., short- and long-horizon tasks). While this perspective effectively measures how well agents execute laboratory tasks, it provides only a partial view of laboratory intelligence. Scientific experiments naturally follow a \textbf{perception--cognition--interaction} loop, requiring agents to continuously perceive laboratory environments, acquire scientific understanding, and perform corresponding experimental operations. Consequently, evaluating agents solely from execution characteristics offers limited insight into the scientific capabilities required for successful experimentation.

This observation motivates a complementary \textbf{capability-oriented} evaluation framework that focuses on the competencies exercised throughout scientific experimentation rather than only execution behavior. Building such a framework presents four key challenges: (1) establishing a principled taxonomy that decomposes laboratory intelligence into representative capability dimensions; (2) enabling scalable generation of diverse laboratory tasks beyond manually curated benchmark suites; (3) constructing realistic simulation environments that faithfully reproduce representative physicochemical processes together with dynamic state transitions; and (4) treating laboratory safety as an inherent component of scientific intelligence rather than an isolated evaluation criterion.

To address these challenges, we present \textbf{SciVLABench}, an open-source benchmark for comprehensive evaluation of scientific embodied intelligence. SciVLABench is built upon four complementary components. First, it establishes a structured evaluation taxonomy that organizes scientific embodied intelligence into representative dimensions spanning \textbf{scientific cognition} and \textbf{scientific practice}. Second, it introduces an automated task generation pipeline that enables scalable construction of large-scale and diverse laboratory tasks beyond manually curated benchmarks. Third, it incorporates a high-fidelity laboratory simulation environment supporting representative physicochemical processes, dynamic laboratory interactions, and event-driven state transitions throughout experimental procedures. Finally, it integrates laboratory safety throughout the evaluation framework, assessing not only task completion but also safety-aware decision-making during experimental execution.

Built upon this framework, SciVLABench currently encompasses more than \textbf{20 representative laboratory scenarios} derived from contemporary scientific literature and automatically generates over \textbf{1,000 task templates} spanning \textbf{six evaluation dimensions}, including chemical knowledge, instrument knowledge, laboratory safety, fundamental laboratory operations, instrument operation, and spatial organization. The benchmark further incorporates more than \textbf{100 specialized laboratory assets}, supports over \textbf{30 robotic manipulation primitives}, covers multiple scientific disciplines including analytical chemistry and biochemistry, and automatically generates expert demonstration trajectories for every task template, providing a scalable resource for both benchmark evaluation and policy learning.

Our contributions are summarized as follows:

\begin{enumerate}
\item We introduce \textbf{SciVLABench}, a comprehensive benchmark that complements existing execution-oriented evaluation with a structured capability-oriented framework for scientific embodied intelligence.

\item We develop an automated task generation pipeline that enables scalable construction of diverse laboratory tasks together with expert demonstration trajectories for both benchmark evaluation and policy learning.

\item We establish a realistic laboratory simulation framework integrating high-fidelity physical modeling, dynamic laboratory interactions, and safety-aware assessment, providing a comprehensive testbed for advancing scientific embodied intelligence.
\end{enumerate}

%%
%% ==================== PLACEHOLDER SECTIONS ====================
\section{Related Work}
\label{sec:related}

\textbf{[TO BE COMPLETED: Related Work Section -- approximately 1.5 pages covering:]}

\begin{itemize}
    \item \textbf{2.1 Embodied AI and Robotics Simulation}: Physics-based simulators (MuJoCo, PyBullet, Isaac Gym), robotics benchmarks (ManiSkill, RoboNet, CALVIN)
    \item \textbf{2.2 Language-Conditioned Manipulation}: VLA models (OpenVLA, RT-series), language grounding, instruction-following benchmarks
    \item \textbf{2.3 Scientific and Lab Automation}: Existing lab automation platforms, scientific experimentation benchmarks, gap analysis
    \item \textbf{2.4 Data-Driven Task Generation}: LLM-based task generation, neuro-symbolic approaches, Real-to-Sim methodologies
\end{itemize}

\textbf{[End of Related Work placeholder]}

\section{The SciVLABench Framework}
\label{sec:framework}

This section presents the technical architecture of SciVLABench, an automated real-to-simulation task generation pipeline based on LangGraph. We describe the overall pipeline design and detail each major stage, from natural-language protocol parsing to physics-based benchmark generation.

\subsection{Overview}

To construct a capability-oriented benchmark, a diverse collection of laboratory task scenarios is required to comprehensively evaluate different scientific capabilities. However, manually designing simulation task templates typically produces only a limited number of isolated scenarios, fundamentally restricting task diversity and making large-scale benchmark construction difficult. To address this limitation, we develop an automated real-to-simulation task generation pipeline based on LangGraph. Given a natural-language description of a laboratory experiment, the pipeline automatically constructs executable simulation tasks through a closed-loop generation–verification–feedback process, enabling scalable, reproducible, and physically validated benchmark generation.

As illustrated in Figure~\ref{fig:framework_overview}, the proposed pipeline consists of three sequential stages that progressively compile real-world laboratory procedures into executable benchmark tasks. Stage I (Environment Construction) compiles a natural-language laboratory protocol into an executable simulation environment by parsing experimental instructions, grounding referenced entities to laboratory assets, and constructing a physically consistent experimental scene. Stage II (Task Decomposition and Planning) compiles the structured protocol into executable robot behaviors by first decomposing the experiment into a sequence of ordered subtasks, and then generating, for each subtask, both an executable manipulation skill sequence and a corresponding success condition for subsequent verification. Stage III (Simulation Verification and Benchmark Generation) validates the synthesized task through physics simulation, verifies task correctness using the generated success conditions, and compiles the verified task into benchmark assets, including executable task templates, expert demonstrations, and trajectory data. To improve robustness and support policy generalization, the validated tasks are further augmented through scene randomization before being incorporated into the benchmark.

From a methodological perspective, the proposed pipeline can be viewed as a progressive task compilation process. Specifically, it successively compiles natural-language experimental protocols into structured laboratory environments, executable manipulation programs, and finally physically validated benchmark instances. This staged compilation framework not only enables scalable and reproducible benchmark construction without manual task engineering, but also tightly couples task generation, physical verification, and expert demonstration collection within a unified automated workflow.

\begin{figure*}[t]
\centering
% TODO: Insert framework architecture diagram
% \includegraphics[width=0.95\textwidth]{figures/framework_overview}
\caption{SciVLABench Pipeline Architecture. The framework consists of three sequential stages: Environment Construction, Task Decomposition and Planning, and Simulation Verification. Each stage progressively compiles natural-language protocols into executable benchmark tasks through a closed-loop generation–verification–feedback process.}
\label{fig:framework_overview}
\end{figure*}

\subsection{Environment Construction}

Stage I compiles a natural-language laboratory protocol into an executable simulation environment by progressively grounding semantic descriptions into physically instantiated laboratory assets. This stage bridges the gap between symbolic laboratory protocols and executable simulation environments through a progressive grounding process. Specifically, it addresses two fundamental challenges: identifying the laboratory entities referenced in the protocol and constructing a physically executable environment that faithfully preserves the semantics of the original experimental procedure. Accordingly, the compilation process consists of two complementary steps: semantic grounding and scene instantiation.

We first perform semantic grounding to identify all laboratory entities together with their semantic attributes, including object categories, quantities, materials, and functional roles in the experiment. Because laboratory protocols frequently involve multiple instances of the same object category (e.g., several beakers or test tubes), purely semantic references become ambiguous during subsequent task planning. We therefore assign each extracted entity a globally unique identifier (UID), establishing unambiguous and persistent symbolic references throughout the entire compilation process. These UIDs remain consistent across task planning, condition generation, and simulation verification, ensuring that every subsequent stage operates on the same physical entity. The grounded entities are subsequently matched against a local laboratory asset repository. If no suitable asset exists, the system automatically retrieves a semantically relevant 3D model from external repositories, allowing the asset library to expand continuously beyond manually curated resources. Retrieved assets are then automatically compiled into simulation-ready laboratory objects by converting them into a unified representation, assigning physically meaningful interaction properties, and validating their geometric consistency, ensuring that every instantiated object can be reliably manipulated during subsequent simulation. As a result, this grounding process establishes a deterministic correspondence between symbolic experimental descriptions and executable physical entities, providing the semantic foundation upon which all subsequent planning and verification stages operate.

After all required assets have been grounded, the simulation environment is instantiated according to the spatial relationships specified in the experimental protocol. Object layouts, relative positions, supporting structures, and workspace organization are automatically generated to preserve both the procedural semantics and physical feasibility of the original experiment. Rather than merely reconstructing the geometric appearance of a laboratory scene, this process compiles a physically executable environment in which every object possesses consistent semantic identities, valid interaction properties, and physically plausible spatial configurations. The compiled environment subsequently serves as the execution substrate for task decomposition, robot skill planning, and success-condition generation in the following stage.

\subsection{Task Decomposition and Planning}

Stage II compiles the executable simulation environment into an executable and verifiable robot program. Rather than directly generating a monolithic manipulation policy from the experimental protocol, the proposed framework progressively synthesizes the robot program through task decomposition and dual-track planning. This design explicitly separates execution from verification, enabling scalable program synthesis while providing reliable evaluation criteria for subsequent physical validation.

We first decompose the laboratory procedure into an ordered sequence of semantically complete experimental steps, each corresponding to an independently executable experimental objective. Instead of splitting the protocol at the sentence level, the proposed strategy groups related manipulations into coherent action clusters that collectively accomplish one meaningful stage of the experiment, such as reagent transfer, solution mixing, or instrument assembly. Beyond simplifying long-horizon planning, this decomposition establishes the fundamental evaluation units of the benchmark, enabling fine-grained process-level assessment rather than binary task-level success or failure. By explicitly aligning planning and evaluation at the step level, the benchmark measures not only whether an experiment is successfully completed, but also how much of the experimental procedure has been correctly accomplished. The resulting ordered experimental steps therefore provide an intermediate representation that bridges high-level laboratory procedures and low-level robotic manipulation.

For each decomposed experimental step, we simultaneously synthesize its execution specification and verification specification through two independent planning processes. The execution branch translates each experimental objective into an executable sequence of atomic manipulation primitives selected from a library of more than 30 reusable robot skills. Instead of relying on manually designed action templates, the language model autonomously composes these primitives according to the semantic objective of each step together with the current environment state, generating a complete manipulation program capable of accomplishing the desired experimental operation. In parallel, an independent verification branch assigns every experimental step an appropriate success condition selected from an expert-defined condition library. These conditions encode the expected physical outcomes of each operation, including object states, spatial relationships, and experimental constraints, allowing every intermediate step to be objectively evaluated during execution. Unlike conventional self-evaluation pipelines, execution synthesis and condition synthesis are performed by two independent language models operating in parallel, preventing the same model from simultaneously generating and evaluating task specifications. More importantly, the two planning branches jointly establish a mutually verifiable program specification. On one hand, the synthesized manipulation program provides a concrete realization of each experimental objective, allowing the generated success conditions to be physically validated. On the other hand, the success conditions serve as explicit execution criteria for verifying whether the synthesized manipulation sequence correctly achieves the intended experimental outcome. This bidirectional consistency substantially improves the reliability of automatically generated benchmark tasks.

Finally, the synthesized atomic manipulation sequences are concatenated according to the decomposed procedure to form a complete expert execution program. Together with the aligned step-wise success conditions, they constitute a fully specified robot program consisting of both executable behaviors and explicit verification criteria. The resulting program is subsequently executed in the physics simulator to automatically produce expert demonstrations and trajectory data, providing both reliable evaluation references for benchmark validation and high-quality supervision for downstream policy learning. Conceptually, Stage II does not merely synthesize expert demonstrations; it automatically constructs a complete program specification that tightly couples execution and verification, laying the foundation for reliable physical validation in the subsequent simulation stage.

\section{The SciVLABench Benchmark}
\label{sec:benchmark}

\textbf{[TO BE COMPLETED: Benchmark Section - approximately 2.5 pages covering task design, evaluation dimensions, data generation, and infrastructure]}

\textbf{Content to include:}
\begin{itemize}
    \item 4.1 Task Design and Coverage
    \item 4.2 Evaluation Dimensions (M\&T, Spatial, CommonSense, Semantic, PhysicalLaw, Complex)
    \item 4.3 Data Generation and Asset Pipeline
    \item 4.4 Evaluation Infrastructure
\end{itemize}

\section{Experiments}
\label{sec:experiments}

\textbf{[TO BE COMPLETED: Experiments Section - approximately 3 pages covering experimental setup, main results, analysis, and ablation studies]}

\textbf{Content to include:}
\begin{itemize}
    \item 5.1 Experimental Setup
    \item 5.2 Main Results (Policy Model Performance, VLM Evaluation)
    \item 5.3 Analysis and Discussion (Physical Verification, Long-Horizon Reasoning, Dynamic Intervention Challenges)
    \item 5.4 Ablation Studies
\end{itemize}

\section{Conclusion}
\label{sec:conclusion}

\textbf{[TO BE COMPLETED: Conclusion and Future Work - approximately 0.5 pages]}

\textbf{Content to include:}
\begin{itemize}
    \item Summary of contributions
    \item Key findings and takeaways
    \item Future directions (expanding task coverage, enhanced intervention modeling, additional robot platforms, real-world validation)
    \item Broader impact statement
\end{itemize}

%%
%% ==================== ACKNOWLEDGMENTS ====================
\begin{acks}
% TODO: Add acknowledgments, funding sources, institutional support
% Example:
% This work was supported by [Funding Agency] under grant [Number].
% We thank [Name] for helpful discussions and [Name] for technical assistance.
\end{acks}

%%
%% ==================== BIBLIOGRAPHY ====================
\bibliographystyle{ACM-Reference-Format}

% TODO: Create references.bib file with your citations
% For now using placeholder - replace 'sample-base' with your bibliography filename
\bibliography{sample-base}

\end{document}